\chapter{Uroboros Scale Quotient and Collatz--Riemann Coordinate Conjugacy}
\label{chap:uroboros-torus}

\status{Exact algebraic coordinate conjugacies proved; torus quotient proved conditional on the stated scale-identification convention. No Collatz-convergence, xi-scaling, zero-location, or RH claim is made.}

\section{Normalized half-layer coordinate}

Let $x>0$ and define
\[
u=2x,\qquad s=\frac{u}{1+u}=\frac{2x}{1+2x}.
\]
Then
\[
\boxed{x=\frac12\iff u=1\iff s=\frac12.}
\]
Thus the distinguished half-layer in the $x$-coordinate is sent exactly to the fixed half-layer of the Riemann reflection coordinate.

Define
\[
I_x(x)=\frac{1}{4x}.
\]
Then
\[
2I_x(x)=\frac1{2x}=\frac1u,
\]
and hence
\[
\frac{1/u}{1+1/u}=\frac1{1+u}=1-\frac{u}{1+u}.
\]
Therefore
\begin{theorem}[Half-centered inversion conjugacy]
Under $s=2x/(1+2x)$,
\[
\boxed{I_x:x\mapsto\frac1{4x}}
\]
is conjugate to
\[
\boxed{J:s\mapsto1-s.}
\]
Equivalently, the reciprocal map $u\mapsto1/u$ is exactly the critical reflection $s\mapsto1-s$.
\end{theorem}

\section{The two Collatz branches in the Riemann coordinate}

For the halving branch $x\mapsto x/2$, direct substitution gives
\[
H(s)=\frac{s}{2-s}.
\]
For the odd branch $x\mapsto3x+1$, direct substitution gives
\[
O(s)=\frac{s+2}{3}.
\]
Thus the standard Collatz map on positive integers is conjugated branchwise to
\[
\boxed{
C_s(s)=
\begin{cases}
\dfrac{s}{2-s}, & x\text{ even},\\[0.75em]
\dfrac{s+2}{3}, & x\text{ odd}.
\end{cases}}
\]
No convergence statement is used in this identity.

The inverse of the halving map is
\[
H^{-1}(s)=\frac{2s}{1+s}.
\]
With $J(s)=1-s$, one verifies
\[
\boxed{JHJ=H^{-1}.}
\]
Hence the halving translation and the half-layer reflection satisfy the standard dihedral conjugation relation.

\section{The fundamental halving chain}

The finite chain
\[
16\to8\to4\to2\to1\to\frac12
\]
contains five exact halving steps and is mapped to
\[
\frac{32}{33}\to\frac{16}{17}\to\frac89\to\frac45\to\frac23\to\frac12.
\]
The endpoint therefore lands exactly on the Riemann half-layer in the $s$-coordinate.

The scale ratio between the first and final $x$-values is
\[
\frac{16}{1/2}=32=2^5.
\]

\section{Scale quotient and complex torus}

Suppose that, as an explicit cycle-identification convention, points related by this scale are identified:
\[
\boxed{u\sim32u.}
\]
This is an added quotient rule, not a consequence of the Collatz conjecture.

For $u\in\mathbb C^*$, write
\[
\lambda=\log u.
\]
Then multiplication by $32$ becomes translation by
\[
\log32=5\log2,
\]
while the multivalued logarithm already identifies
\[
\lambda\sim\lambda+2\pi i.
\]
Therefore the quotient lattice is
\[
\boxed{\Lambda=5\log2\,\mathbb Z+2\pi i\,\mathbb Z.}
\]
Consequently
\begin{theorem}[Uroboros scale quotient]
Under the explicit identification $u\sim32u$,
\[
\boxed{\mathbb C^*/32^{\mathbb Z}\cong\mathbb C/\Lambda}
\]
is a complex one-dimensional torus.
\end{theorem}

\section{INTERPRETACJA: Uroboros cycle semantics}

\textbf{INTERPRETACJA.} The quotient identification $u\sim32u$ can be read as a mathematical realization of the proposed Uroboros convention: the end of the selected five-halving cycle is glued to the beginning of the next cycle. In this interpretation, the half-layer is not a terminal state but a cycle seam.

This interpretation is deliberately separated from the proved algebraic conjugacies. In particular, it does not assert that every Collatz orbit enters this chain.

\section{Relation to the Riemann xi function}

For
\[
X(u)=\xi\!\left(\frac{u}{1+u}\right),
\]
the xi functional equation supplies the already exact inversion symmetry
\[
\boxed{X(u)=X(1/u).}
\]
The present construction supplies a separate discrete scale quotient $u\sim32u$ only after the explicit Uroboros identification convention is imposed.

No identity
\[
X(32u)=X(u)
\]
or quasi-periodic law
\[
X(32u)=M(u)X(u)
\]
has been proved. Such a scaling or $q$-difference law, if independently derived rather than assumed, would constitute a separate analytic statement connecting the torus dynamics and the xi function.

\section{Proof firewall}

The following remain OPEN:
\begin{itemize}
\item convergence of all Collatz trajectories;
\item a natural xi scaling or quasi-periodicity under $u\mapsto32u$;
\item a zero-preserving consequence of the torus quotient;
\item SOH-G003 real-rootedness;
\item the Riemann Hypothesis.
\end{itemize}
